\subsubsection{\href{https://gitlab.ti.bfh.ch/groups/gausf1-vonal3/-/milestones/19}{Research: AI Integration}}\label{subsubsec:milestone-research-ai-integration}
\textit{10.03.2026 -- 22.03.2026}

The goal of this research milestone was to investigate how \glsxtrshort{AI} can support automated security testing and adversary emulation.
The research evaluated potential integration points where \glsxtrshort{AI} could enhance scenario execution, automation, decision-making, and analysis capabilities of the \glsxtrshort{NSAK} framework.
The milestone was extended by three days following input from the expert and instructor to clarify the intended interaction between \glsxtrshort{LLM}s and \glsxtrshort{NSAK}.

\begin{itemize}
    \item[$\square$] How is \glsxtrshort{AI} currently used in cybersecurity automation?
    \item[$\square$] Which tasks in adversary emulation can be supported or automated by \glsxtrshort{AI}?
    \item[$\square$] How could \glsxtrshort{AI} interact with the \glsxtrshort{NSAK} scenarios and drills (e.g.\ Agents, \glsxtrshort{MCP})?
    \item[$\square$] What architectural changes would be required to integrate \glsxtrshort{AI} into \glsxtrshort{NSAK}?
\end{itemize}
